\section{Introduction}\label{sec:v8-introduction}

Political manifestos are long measurement objects. A party platform can run
for dozens or hundreds of pages, and the target score is a document-level
expert judgment about a policy dimension. Economic policy, for example, is
spread across fiscal pledges, welfare commitments, industrial strategy,
privatization language, and qualifications about when each commitment applies.
The measurement target is global. The evidence is local and dispersed.

Large language models make this tension more visible. They can read and score
long texts, but high predictive accuracy on LLM-produced surrogate labels does
not by itself protect downstream empirical estimates from bias
\citep{EgamiEtAl2024}. Long inputs add another failure mode: relevant evidence
can be missed or diluted when it is buried in the middle of a context window
\citep{LiuEtAl2024LostMiddle}. A single full-context score gives the analyst
one number. It does not say which policy distinctions were preserved, which
were compressed away, or where an expert should audit the result.

The Manifesto Project illustrates the cost side of the problem. Trained coders
and expert surveys provide trusted political-text measurements, and the
Manifesto corpus supplies party election programs across countries and time
\citep{ManifestoProject2025a,ManifestoCorpus2025}. Each new full-document
expert target is expensive because the unit is the whole platform. The same
corpus also contains smaller evidence units, especially quasi-sentence coding,
that suggest a different supervision design: audit and train at the span or
merge level when those units measure the same policy target.

C-TreePO makes that design explicit. A compression tree splits a manifesto into
contiguous leaves, summarizes each leaf into a policy-evidence state, merges
sibling states up the tree, and scores the root. The learned state map is
valid for a target oracle \(f^\ast\) only when it satisfies three local laws:
leaf summaries preserve the relevant evidence in their spans, produced states
can be reused without policy drift, and merged states preserve the represented
parent span. These laws are local, but their consequence is global: a valid
root state can replace the original manifesto for oracle-compatible readouts.

The objective has the same two-level shape. The root channel fits the observed
document-level expert target. The local-law channel trains and audits the
state map inside the realized tree. In LLM practice, the dense local signal is
usually a learned proxy \(\widehat f\) or a teacher trace, while direct
\(f^\ast\) calls are sparse. C-TreePO therefore uses \(\widehat f\) for local
training geometry and sampled node-level \(f^\ast\) calls, with logged
propensities, to correct proxy bias in the local-law channel.

The empirical application is the six-dimension Manifesto/Benoit benchmark
\citep{BenoitEtAl2025}. The current evidence is root-observed: held-out
platforms have external expert-survey means, and the C-Tree root is scored
against those document targets. The same runs produce readable leaves, merges,
and proxy diagnostics that define the next audit layer. That separation is the
point of the framework. Root accuracy measures the final task; node audit
measures whether the final score is carried by locally faithful policy states.

Section~\ref{sec:v8-manifesto-mechanism} introduces the manifesto policy-state
mechanism. Section~\ref{sec:v8-ctree-math} states C-Trees and the local laws.
Section~\ref{sec:v8-objective} gives the objective and theorem ladder.
Section~\ref{sec:v8-audit-certificate} gives the audit estimator and
finite-sample certificate. Section~\ref{sec:v8-manifesto-results} reports the
Manifesto application. Section~\ref{sec:v8-scope-related} places the work
against measurement, long-context, compression, sketching, preference-learning,
and neural-operator literatures. Section~\ref{sec:v8-discussion} discusses
how to read the certificate, spend root and node labels, and diagnose
manifesto failure modes.
